% Problem-section file following ../_template.tex.
\section{Capacity-achieving codes for amplitude damping}
\label{sec:problem-2}

\paragraph{Problem.}
What is the constructive quantum code for achieving the quantum capacity
$\mathcal{Q}(\mathcal A_p)$ of the qubit amplitude-damping channel
\begin{equation}
  \mathcal A_p(\rho)=A_0\rho A_0^\dagger+A_1\rho A_1^\dagger,
  \qquad 0\le p\le1?
  \label{eq:p2-amplitude-damping}
\end{equation}
The operators appearing in Eq.~\eqref{eq:p2-amplitude-damping} are
\begin{equation}
  \begin{aligned}
    A_0&=\lvert0\rangle\!\langle0\rvert
         +\sqrt{1-p}\,\lvert1\rangle\!\langle1\rvert
       =\begin{pmatrix}1&0\\0&\sqrt{1-p}\end{pmatrix},\\
    A_1&=\sqrt p\,\lvert0\rangle\!\langle1\rvert
       =\begin{pmatrix}0&\sqrt p\\0&0\end{pmatrix}.
  \end{aligned}
  \label{eq:p2-kraus-operators}
\end{equation}
Equation~\eqref{eq:p2-kraus-operators} uses $p$ as the decay probability of
the excited state.

\subsection*{Status}

Solved

\subsection*{Source}

This constructive gap is implicit in the exact capacity formula of Wolf and
P\'erez-Garc\'ia and the polar-code construction of Wilde and Guha, which
attains only the symmetric coherent-information rate for this channel
\sourcecite{ref:p2-wolf}{WPG07},
\sourcecite{ref:p2-wilde-guha}{WG13}.

\subsection*{Progress}

\begin{itemize}
  \item The channel is degradable for $0\le p\le1/2$ and antidegradable for
  $1/2\le p\le1$.  Consequently,
  \begin{equation}
    \mathcal{Q}(\mathcal A_p)=
    \begin{cases}
      \displaystyle\max_{0\le q\le1}
      \bigl[h_2((1-p)q)-h_2(pq)\bigr],&0\le p\le\tfrac12,\\[1.5mm]
      0,&\tfrac12\le p\le1,
    \end{cases}
    \label{eq:p2-capacity}
  \end{equation}
  where $q$ is the excited-state population of the diagonal input and
  $h_2(x):=-x\log_2x-(1-x)\log_2(1-x)$, with $0\log_2 0:=0$.  The capacity formula in
  Eq.~\eqref{eq:p2-capacity} follows from the small-environment analysis
  \sourcecite{ref:p2-wolf}{WPG07}.

  \item For $0\le p\le1/2$, Wilde and Guha constructed a channel-adapted
  quantum polar code for degradable channels with a classical environment,
  explicitly including the amplitude-damping channel.  It achieves
  \begin{equation}
    I_c(I/2,\mathcal A_p)
    =h_2\!\left(\frac{1-p}{2}\right)-h_2\!\left(\frac p2\right),
    \label{eq:p2-symmetric-rate}
  \end{equation}
  the symmetric coherent-information rate.  Equation~\eqref{eq:p2-symmetric-rate}
  is the bracket in Eq.~\eqref{eq:p2-capacity} evaluated at $q=1/2$; since the
  maximizer of Eq.~\eqref{eq:p2-capacity} is interior and $q=1/2$ is not a
  stationary point of the bracket for $0<p<1/2$, the symmetric rate is
  strictly smaller than $\mathcal{Q}(\mathcal A_p)$ throughout that range.
  The rate in Eq.~\eqref{eq:p2-symmetric-rate} is obtained with an efficient
  encoder and
  asymptotically vanishing entanglement consumption, although that work did
  not establish an efficient decoder \sourcecite{ref:p2-wilde-guha}{WG13}.
  Wilde and Renes subsequently gave a polar construction for arbitrary
  qubit-input channels that achieves the same rate using coherent
  successive-cancellation decoding \sourcecite{ref:p2-wilde-renes}{WR12}.

  \item Renes closed the remaining asymmetric-input and explicit-decoder
  gaps.  In the CSS-type quantum polar scheme, the amplitude part of
  $\mathcal A_p$ is a classical $Z$-channel whose capacity-achieving input
  is nonuniform, so the amplitude code is a Honda--Yamamoto polar code for
  asymmetric channels \sourcecite{ref:p2-honda-yamamoto}{HY13}, while the
  phase part is a pure-state channel decoded by a quantum belief-propagation
  circuit.  Renes states that the rate of this construction, optimized over
  the amplitude input distribution, equals the capacity in
  Eq.~\eqref{eq:p2-capacity}, that degradability of $\mathcal A_p$ removes
  the need for preshared entanglement, and that the decoder is an explicit
  circuit with $O(N^2)$ gates for block length $N$
  \sourcecite{ref:p2-renes-belief-propagation}{Ren17}.
\end{itemize}

\subsection*{References}

\begin{enumerate}[label={},leftmargin=4.8em,labelsep=0.5em]
  \footnotesize
  \item[\textup{[WPG07]}]\label{ref:p2-wolf}
  M. M. Wolf and D. P\'erez-Garc\'ia, ``Quantum Capacities of Channels with
  Small Environment,'' \emph{Physical Review A} \textbf{75}, 012303 (2007).
  \href{https://doi.org/10.1103/PhysRevA.75.012303}{doi:10.1103/PhysRevA.75.012303};
  \href{https://arxiv.org/abs/quant-ph/0607070}{arXiv:quant-ph/0607070}.

  \item[\textup{[WG13]}]\label{ref:p2-wilde-guha}
  M. M. Wilde and S. Guha, ``Polar Codes for Degradable Quantum Channels,''
  \emph{IEEE Transactions on Information Theory} \textbf{59}, 4718--4729
  (2013). \href{https://doi.org/10.1109/TIT.2013.2250575}{doi:10.1109/TIT.2013.2250575};
  \href{https://arxiv.org/abs/1109.5346}{arXiv:1109.5346}.

  \item[\textup{[WR12]}]\label{ref:p2-wilde-renes}
  M. M. Wilde and J. M. Renes, ``Quantum Polar Codes for Arbitrary Channels,''
  in \emph{2012 IEEE International Symposium on Information Theory
  Proceedings}, 334--338 (2012).
  \href{https://doi.org/10.1109/ISIT.2012.6284203}{doi:10.1109/ISIT.2012.6284203};
  \href{https://arxiv.org/abs/1201.2906}{arXiv:1201.2906}.

  \item[\textup{[HY13]}]\label{ref:p2-honda-yamamoto}
  J. Honda and H. Yamamoto, ``Polar Coding Without Alphabet Extension for
  Asymmetric Models,'' \emph{IEEE Transactions on Information Theory}
  \textbf{59}, 7829--7838 (2013).
  \href{https://doi.org/10.1109/TIT.2013.2282305}{doi:10.1109/TIT.2013.2282305}.

  \item[\textup{[Ren17]}]\label{ref:p2-renes-belief-propagation}
  J. M. Renes, ``Belief Propagation Decoding of Quantum Channels by Passing
  Quantum Messages,'' \emph{New Journal of Physics} \textbf{19}, 072001
  (2017). \href{https://doi.org/10.1088/1367-2630/aa7c78}%
       {doi:10.1088/1367-2630/aa7c78};
  \href{https://arxiv.org/abs/1607.04833}{arXiv:1607.04833}.
\end{enumerate}

\subsection*{Comment}

Because the amplitude-damping channel is nonunital, the maximizing population
$q$ in Eq.~\eqref{eq:p2-capacity} is generally not $1/2$, so the symmetric
coherent-information rate in Eq.~\eqref{eq:p2-symmetric-rate} is not the
quantum capacity, and codes attaining only that rate do not settle the
problem.  The solved status records Renes' theorem that the asymmetric polar
construction attains the optimized rate in Eq.~\eqref{eq:p2-capacity} with
an explicit capacity-achieving decoder, thereby solving the problem in the
asymptotic coding sense.  Reducing the decoder complexity below $O(N^2)$
remains a separate algorithmic question.

\subsection*{Tag}

Amplitude-damping channels; Quantum capacity; Quantum coding theory;
Quantum polar codes; Degradable channels; Qubit systems

\subsection*{ID}

\texttt{op\_fcd21a1a5021e464}
